\chapter{Glossary}
\label{app:glossary}
One-line definitions for quick lookup. Each term is developed properly in the chapter
that introduces it; the formula index (\cref{app:formulas}) says where.

\begin{description}
\item[Admittance] Reciprocal impedance, \(Y=1/Z\), measured in siemens.
\item[Bandwidth] Frequency span over which a circuit meets a stated response criterion.
\item[Bode plot] Magnitude and phase of a frequency response versus logarithmic frequency.
\item[Characteristic impedance] Traveling-wave voltage/current ratio on a transmission line.
\item[Corner frequency] Frequency where a first-order response begins its asymptotic slope change.
\item[Cutoff frequency] Often the \(-\SI{3}{\decibel}\) frequency for a filter.
\item[Damping ratio] Dimensionless \(\zeta\) grading a second-order system; \(\zeta=1/(2Q)\).
\item[Gain margin] How far below \(\SI{0}{\decibel}\) the loop gain sits where its phase reaches \(-180^\circ\).
\item[Gain--bandwidth product] Constant product of closed-loop gain and bandwidth for a single-pole amplifier.
\item[Barkhausen criterion] Oscillation condition \(L=-1\): \(|L|=1\) with \(\angle L=-180^\circ\).
\item[Impedance] Complex ratio \(Z=V/I\) for sinusoidal steady state.
\item[Loaded Q] Quality factor including external loading or coupling.
\item[Loop gain] Gain around an opened feedback loop, \(L=A\beta\); sets stability and how insensitive the closed loop is to changes in the forward path. (Not to be confused with receiver \emph{desensitization}, an unrelated exam term meaning loss of sensitivity caused by a strong nearby signal.)
\item[Phase margin] Amount by which loop phase exceeds \(-180^\circ\) at unity loop gain.
\item[Phasor] Complex representation of sinusoidal magnitude and phase.
\item[Pole] Root of a transfer-function denominator; equivalently an eigenvalue and a natural mode.
\item[Reactance] Imaginary part of impedance associated with energy storage.
\item[Root locus] Path traced by the poles as a parameter (e.g.\ \(R\) or feedback) varies.
\item[Quality factor \(Q\)] Sharpness of a resonance, \(Q=1/(2\zeta)=2\pi\times\) (energy stored)/(energy lost per cycle). \emph{Component} \(Q\) is one part's loss; see also Loaded \(Q\), Unloaded \(Q\).
\item[Reflection coefficient] Complex ratio of reflected to incident wave amplitude.
\item[Return loss] Decibel measure of reflection magnitude.
\item[Roll-off] Rate at which a response falls past its passband, quoted in decibels per decade or octave.
\item[Ringing] The decaying oscillation of an underdamped second-order system.
\item[SWR] Standing-wave ratio, \((1+|\Gamma|)/(1-|\Gamma|)\); 1 means a perfect match.
\item[State variable] Variable that captures stored energy needed to predict future behavior.
\item[Susceptance] Imaginary part of admittance, \(B\); for a pure reactance \(B=-1/X\).
\item[Unloaded Q] Quality factor due to internal resonator loss alone.
\item[Velocity factor] \(VF=v_p/c\), how much slower a wave travels in a line than in free space.
\item[Virtual ground] A node held near zero volts by feedback rather than by a physical connection.
\item[Zero] Root of a transfer-function numerator; a frequency of null transmission (e.g.\ a trap).
\end{description}
